\section{Проблемы, ограничения и перспективы развития МКВ}
\label{sec:limitations}

Вычислительные и коммуникационные издержки, скрытый доступ к данным, доверие к среде исполнения, утечки через результаты и организационные барьеры определяют ограничения внедрения МКВ и выбор решений с учётом производительности и модели безопасности.

\subsection{Коммуникационные издержки и гибридные протоколы}
\label{subsec:communication}

Защищённое вычисление требует дополнительных криптографических операций и обменов. При разделении секрета линейные операции часто выполняются локально, а умножение требует взаимодействия. В зашифрованных схемах трафик создаёт передача таблиц нелинейных логических элементов. Рост схемы увеличивает объём обмена, потребление памяти и время выполнения; при высокой сетевой задержке ограничением становится число последовательных раундов \cite{06}.

Для снижения издержек гибридные протоколы сочетают разделение секрета, зашифрованные схемы и гомоморфное шифрование: арифметические операции можно выполнять над секретными долями, сравнения --- булевыми схемами, подходящие групповые операции --- над шифротекстами. Представление выбирают по структуре задачи и стоимости переходов между примитивами; их объединение не гарантирует ускорения. Дополнительно применяют пакетную обработку и предварительную генерацию вспомогательных криптографических данных \cite{06, 08}.

В МО коммуникационные расходы дополняются трудностями представления дробных чисел, усечения и нелинейных функций. Гибридный подход требует выбора фиксированной точности и оптимизации свёрток и функций активации. В схеме разделения секрета Шамира \cite{12} ускорение оценивается совместно с точностью модели и устойчивостью к переполнению (см. подраздел~\ref{subsec:ml-applications}).

Автоматизированный выбор криптографического представления операций должен учитывать число участников, сетевые задержки и ресурсы, сохраняя заданную модель противника. Сокращение трафика не устраняет преждевременное прекращение протокола: безопасность с прерыванием, справедливость получения результата и устойчивость к отказам остаются самостоятельными требованиями \cite{08}.

\subsection{ORAM, масштабируемость и управляемые утечки}

Адреса и последовательность запросов могут раскрывать сведения об операции даже при защите содержимого записей. Для сокрытия шаблонов доступа применяются протоколы ORAM (oblivious random access memory). Дополнительные обращения, обмены и расход памяти усложняют масштабирование МКВ \cite{06}.

Накладные расходы можно снизить, допуская управляемую утечку заданных характеристик: размера набора или ограниченных сведений о структуре обращений при сохранении защиты остальных данных. Допустимую утечку необходимо явно включить в модель безопасности и анализировать с учётом дополнительной информации противника.

Выбор между ORAM и управляемыми утечками зависит от чувствительности шаблона обращений. Если его раскрытие недопустимо, издержки снижают организацией данных, пакетной обработкой и оптимизацией скрытого доступа. При допустимом раскрытии упрощение доступа может уменьшить стоимость вычисления. Повторные запросы требуют учёта накопления информации: приемлемость одного раскрытия не гарантирует безопасности последовательности.

Специализированные структуры данных и протоколы согласуют эффективность с заданным профилем утечки. Их оценивают по времени выполнения, памяти и составу раскрываемых сведений; компромисс производительности и утечки не заменяет формального обоснования конфиденциальности.

\subsection{Доверие к среде исполнения и TEE}

Доказательство безопасности протокола не исключает ошибок реализации, компрометации устройств, уязвимостей программно-аппаратного стека обработки долей и побочных каналов утечки. Защита от активного противника не отменяет требований к безопасности устройств и ключевого материала честных участников \cite{06}.

Для ограничения доступа к чувствительным данным применяются доверенные среды выполнения (TEE). Они позволяют изолировать выбранные операции от остального программного окружения. В гибридной архитектуре TEE выполняет часть обработки, а МКВ обеспечивает распределённое вычисление. Подход может снижать стоимость операций, но безопасность дополнительно зависит от аппаратной платформы и механизмов её защиты.

Загартдинов и Поляков при сравнении технологий конфиденциальных вычислений подчёркивают необходимость учитывать жизненный цикл защищаемого программного обеспечения и сложности интеграции аппаратных механизмов с существующими средствами защиты \cite{01}. Румянцев и Хаджиева рассматривают защищённую вычислительную среду через взаимосвязь физической и логической составляющих защиты и надёжности вычислений \cite{03}.

В настоящей работе из сопоставления этих положений сделан вывод: TEE не устраняет все риски, а её гарантии не следует автоматически приравнивать к криптографическим гарантиям МКВ. Проектирование требует учёта доверия к производителю, проверки состояния кода, управления ключами, обновления компонентов и защиты от побочных каналов. TEE дополняет гибридную архитектуру, но не заменяет МКВ без дополнительных условий.

Направления разработки включают проверяемое развёртывание, аттестацию среды и сочетание аппаратной изоляции с распределением секретов. Архитектура должна определять доступные внутри TEE данные, гарантии при её компрометации, процедуры обновления и восстановления. Применение ускорения требует приемлемой для участников модели доверия.

\subsection{Утечка через результаты и дифференциальная приватность}

МКВ защищают входные данные в процессе вычисления, однако не ограничивают информацию, которую раскрывает разрешённый результат. Точное значение агрегата, предсказание модели или последовательность ответов могут позволить сделать выводы о данных отдельного участника. Такая утечка не обязательно означает нарушение протокола: она может следовать из самой вычисляемой функции \cite{06}.

Для ограничения раскрытия через выход МКВ сочетают с дифференциальной приватностью. Результат формируют с контролируемой случайной составляющей, параметры которой выбирают с учётом чувствительности запроса и заданных требований к приватности. Если точный результат нельзя раскрывать ни одной стороне, добавление шума должно быть частью защищённого вычисления до выдачи ответа. Таким образом, МКВ защищают процесс обработки, а дифференциальная приватность ограничивает сведения, доступные из опубликованного результата \cite{06, 13}.

В распределённом МО защищённое вычисление с дифференциальной приватностью можно применять к агрегированию обновлений или публикации статистик. Федеративное обучение позволяет оставить исходные данные у их владельцев, но само по себе не устраняет утечки через передаваемые обновления; локальное обучение, защищённое агрегирование и дифференциальная приватность выполняют разные функции, согласуемые в общей архитектуре \cite{13}.

Шум может снижать точность модели, а повторные запросы требуют учёта совокупного раскрытия. Компромисс приватности и полезности оценивают через параметры шума и бюджет приватности совместно с точностью, временем выполнения и сетевым трафиком. Необоснованное добавление шума не обеспечивает заявленных гарантий.

\subsection{Правовые и организационные барьеры и стандартизация}
\label{subsec:standardization}

Внедрение МКВ требует согласования обязанностей участников, условий обработки и использования результатов. Криптографическая защита не заменяет правового основания обработки и определения ответственности за ошибки, прекращение вычисления или раскрытие информации. Различия терминологии, моделей безопасности и интерфейсов, а также потребность в специалистах для сопровождения системы на протяжении жизненного цикла создают дополнительные трудности \cite{01, 13}.

Стандартизация требует единообразного описания модели противника, допустимого сговора, правил выдачи результата, управляемых утечек и зависимостей от доверенных компонентов. Общие интерфейсы и методики испытаний позволяют сопоставлять совместимость, производительность и безопасность при согласованных предположениях.

Организационные процедуры фиксируют права доступа, подключение участников, обновление ключей и компонентов, действия при отказах и публикацию результатов. Проверка реализации и документации дополняет выбор примитива. Стандартизация не гарантирует правового соответствия: применимость решения оценивается для конкретной задачи и условий обработки.

Направления разработки включают согласование терминологии, профили безопасности типовых приложений и воспроизводимые испытания. Обновление криптографических механизмов, в том числе при изменении модели угроз, предусматривается в жизненном цикле системы. Эти процедуры связывают прототипы с проверяемыми правилами промышленного внедрения; наличие стандарта не заменяет доказательства безопасности.

\subsection{Вывод}

Подходы, рассмотренные в подразделах~\ref{subsec:communication}--\ref{subsec:standardization}, изменяют разные параметры системы: издержки вычисления и скрытого доступа, доверие к среде, раскрытие результатов и сопоставимость требований внедрения.

Ни один подход не устраняет все ограничения; архитектуру оценивают по производительности, точности, допустимому раскрытию и модели доверия. В следующем разделе эти критерии применены к проекту защищённого МО.